% =====================================================================================================
% Discussion (the \section and the label sec:discussion are in main.tex; this file has no \section).
% Every number is quoted from Section sec:experiments (same sources: results/digest.md, results/*_summary.json,
% results/findings.md). The A4 and A4b sweeps are final (results/A4_summary.json, results/A4b_summary.json,
% findings.md Sec. 7): slopes are least-squares fits of the per-kappa geometric means over kappa >= 10n, with 95%
% t-intervals with K-2 degrees of freedom. The hypotheses are stated in the introduction (paragraph Hypotheses).
% Citations: claims as in related.tex and bibliography/literature_review.md (Sections 6, 8, 9); for A1, the
% assessment of hypothesis (A1) in Appendix A1 (theory/A1_lower_bound.tex, Sec. lb:statements, items (i)-(iv)).
% A3: A3 note, Outcome and the reading of q and q_+ in Sec. lk-instrument (A3_summary.json: q >= 1 on 63%, 5%, 65% of
% the steps; q_plus_frac_ge_1 0.0986, 0.0519, 0.1929 for (H), logistic kappa = 1e3, 1e4). Growth from kappa = 1e3 to
% 1e4 (alpha = 1): 4.117/2.775 = 1.48 (A3_summary.json, speedup_exact_over_lemma) against 15.734/5.155 = 3.05
% (B3_summary.json, attribution, speedup).
% Results of the appendices are referred to with \resref{label} (report/restate_defs.tex).
% B1: B1 note, Step 2 of the proof and outcome paragraph (E V_t^P >= 3 gamma^2 (1-1/n)^t M_0); B1.pkl (memory
% coefficient about 1-0.65/n, r_sq 0.51-0.53, r_var about 1). A2: A2_summary.json and A2.pkl (with storage of the
% reference gradients, i.e. grad_evals - iters: L-Katyusha 71400 against SAGA 87800 on logistic regression).
% A4: naive/L-Katyusha at kappa = 1e4 (A4_summary.json): 3.01 on (H), 2.99 on (R).
% =====================================================================================================

\paragraph{B3: where the proof of Proposition~\ref{prop:59} loses.}
The diagnosis was carried out and locates the loss, which depends on the regime. On $\ell_2$-regularized logistic
regression ($n=1000$, $d=50$, 5 seeds) the exact one-step contraction exponent of SAGA, averaged over passes 2--20,
exceeds the certified one by a factor $3.2$--$3.7$ for $\kappa\le10^2$ and $19$--$220$ at $\kappa=10^5$. For $\kappa\gg
n$ the steps that cost are those that replace a quantity depending on the iterate by a global constant: for
$\kappa\ge10^4$, cocoercivity with $R^2=\max_iL_i$ (I2) and strong convexity with $\mu=\lambda$ (I4) carry
$0.77$--$0.96$ of the logarithmic gap in proof order. For $\kappa\le10^2$, where the memory term is active, the gap is
carried in proof order by the square-norm step I1 ($0.39$--$0.44$) and by I2 ($0.32$--$0.37$), and in order~A by I5
and I1. The variance step (I3) is nearly tight ($r_{\rm var}\ge0.958$; at most $0.6\%$ of the gap in proof order,
$6.3\%$ in order~A) and the minimum of the two rates (I5) matters only when $\kappa\lesssim n$. The additive shares hide
a multiplicative structure: the $\kappa$-dependent part of the gap is the factor $\muEff/\mu$ of I4, whereas, by
\resref{prop:B3-representation}, I1 and I2 together can contribute at most $\frac83$ at the fixed step (measured
$2.45$--$2.62$), and the product of the two reproduces $91$--$96\%$ of the speed-ups for $\kappa\ge10^4$. By
\resref{prop:B3-optimum} and~\resref*{rem:B3-conclusion}, optimizing $(\gamma,\Delta)$ within the proof gains $17\%$
at $\kappa=10^3$ (at most $\frac43$ whenever $\kappa\ge\frac{9n}{16}$), but up to a factor $2.4$--$2.9$ for
$\kappa\le10^2$, comparable to the whole gap there; and, by \resref{prop:B3-lower}, no analysis of SAGA with a constant
step can remove the exponent $1$ of $\kappa$. The theorem has the right order in the worst case; on these instances it
is pessimistic by a factor $3.2$--$3.7$ already for $\kappa\le10^2$, and its pessimism grows with $\kappa$ because the
curvature met along the trajectory exceeds $\mu$ by a factor that grows with $\kappa$ (up to $81$ here).

\paragraph{B2: trajectory-dependent constants.}
Confirmed in its main claim, with one refuted prediction. By \resref{cor:B2-local}, near $\etastar$, $\muEff/\mu$ is
bounded below by about $\lmin(F''(\etastar))/\mu$ (up to $70$ in our problems), a value that SAGA reaches when its error
aligns with the flattest direction, and $R^2$ overestimates $\LEff$ by a factor $2.6$ to $86$; the a-posteriori bound
with $\muEff$ and $\LEff$ (\resref*{prop:B3-traj}) recovers $0.59$--$0.96$ of the logarithmic gap for $\kappa\ge10^4$,
and nothing for $\kappa\le10^2$, where the bound is saturated by the memory term. Our prediction that the error of SAGA
stays spread over many directions holds only for $\kappa\lesssim n$, and not in every run with $\kappa\approx n$: at
$\kappa=10^4=10n$ and for $\alpha\in\{0,1\}$ (one run each) SAGA keeps and sharpens, as gradient descent does, the
alignment of the error with the flattest direction of $F''(\etastar)$, on which the initial error already carries
$89$--$98\%$ of its squared norm, a behaviour that \resref{rem:B2-P2-outcome} explains only heuristically; for
$\alpha=2$, whose spectrum is nearly degenerate at the bottom, neither method aligns within our runs. The effective constants are
a-posteriori quantities: they explain an observed trajectory, they do not certify one in advance, and, by
\resref{rem:B2-what-explains}, at a fixed step the gain from $\LEff$ is capped at $\frac83$.

\paragraph{B1: preconditioned SAGA.}
Nuanced. For logistic regression no fixed metric improves the global condition number: $\kappaP\ge\kappa$ by
\resref{prop:B1-nofreelunch}, so that \resref{cor:B1-rate} never favours preconditioning. Yet at $\kappa=10^3$ the
statistical and Hessian metrics raise the observed exponent per pass (in component gradients; a dense metric costs
$O(d^2)$ instead of $O(d)$ per iteration, and $P=F''(\etastar)$ requires $\etastar$, an oracle benchmark) from $1.00$
to $1.11$ ($\alpha=1$) and from $0.61$ to $1.24$--$1.26$ ($\alpha=2$). The gain is local: the trajectory bound of
\resref{cor:B1-kappa-traj}, which pays the smallest per-sample curvature, does not explain it, whereas with the local
constant at the optimum, a heuristic substitution that \resref{lem:B1-asymptotic} supports only as
$\theta\to\etastar$, the bound reproduces the observed ordering but saturates at the memory term, $\frac13$ per pass,
$3.3$--$3.8$ times below the observed exponents. Along the trajectories, the exact memory coefficient accounts for a
factor of about $2$ of this gap (about
$1-\frac{0.65}n$ instead of $1-\frac1{3n}$), and the parameter part, which contracts faster but is charged at the memory
rate by I5, for most of the rest; but no bound on $\E V_t^P$ can gain more than a factor $3$ over $\frac13$, since the
memory slots that are never refreshed keep $\E V_t^P$ from decreasing faster than $(1-\frac1n)^n\approx e^{-1}$ per
pass: a factor $1.1$--$1.26$ separates the typical trajectories, whose slopes are measured, from the expectation that
the bound controls. The metric shifts the slack of the proof from cocoercivity to strong convexity, as predicted, and
the Jacobi metric does not help when the features are rotated.

\paragraph{A1: lower bounds.}
This hypothesis is theoretical; Appendix~\ref{sec:lower_bound} (Section~\ref{sec:th-A1}) assesses it against the known
bounds, restated with their exact hypotheses. The first clause holds on instances whose dimension is of the order of
the number of queries certified: span-based methods sampling i.i.d.\ from a fixed distribution need
$\Omega((n+\sqrt{n\kappa})\log\frac1\varepsilon)$, by \resref{thm:LZ} when $\kappa\ge n$ and, on the instance (S), for
every $\kappa>1$ by \resref{cor:count}. The second clause holds only in a restricted sense: for oblivious methods (up
to logarithmic factors) and for loopless methods with a full-gradient coin, full-gradient steps can divide the $n\log\frac1\varepsilon$ term by at most $1+(\log(n/\kappa))_+$, a gain only
when $\kappa\ll n$, and leave the $\sqrt{n\kappa}\log\frac1\varepsilon$ term untouched, by part~(c) of
\resref{thm:HLOY} and by \resref{thm:HXZ}. Deterministic methods, cyclic ones included, face instead
$\Omega(n\sqrt\kappa\log\frac1\varepsilon)$ for small enough $\varepsilon$ when $\kappa>73$, by \resref{thm:WS}. Two
questions remain open: for randomized methods with fully adaptive queries only an additive $n$ is proved (randomized
part of \resref*{thm:WS}), and whether randomly reshuffled span-based methods, or deterministic ones when $\kappa\le73$,
can beat $n\log\frac1\varepsilon$ is not settled. Our experiments use the lower bound only through its instance: the
exponent $\frac12$ is certified there inside the window of $\kappa$ of \resref{prop:A4}, which the dimension limits
($\kappa_{\max}\approx4.9\cdot10^4$ for $m=256$ and $\varepsilon=10^{-8}$), and the regime $\kappa\ll n$ of the second
clause was not explored numerically.

\paragraph{A3: corrected acceleration.}
Confirmed. The one-step contraction of the Lyapunov proof of L-Katyusha, \resref{thm:lk}, holds at every recorded
state and is pessimistic in exponent by a factor $2.8$--$4.7$ on our three problems ($3.7$--$9.3$ for its closed-form
rate). On the two logistic problems this is less than for Proposition~\ref{prop:59} ($2.8$ and $4.1$ against $5.2$ and
$15.7$) and grows less from $\kappa=10^3$ to $10^4$ ($\times1.5$ against $\times3.1$); two values of $\kappa$, in
different regimes of the certificate, do not show how it scales. Its variance bound loses, as I2 does, the gap between
the worst-case curvature $\max_iL_i$ and the curvature met along the trajectory (tightness $0.125$--$0.15$). The term
that the negative momentum cancels in \resref{prop:cancel} exceeds the whole certified decrease in absolute value on
$63\%$ and $65\%$ of the steps of the two ill-conditioned problems, but mostly with the sign that helps a bound without
the anchor: it is positive and exceeds this decrease on $9.9\%$ and $19.3\%$ of the steps ($5.2\%$ on logistic
regression with $\kappa=10^3$), and only there could it not have been absorbed, within this proof, into the certified
decrease instead of cancelled.

\paragraph{A2: naive acceleration.}
The failure of naive acceleration is confirmed on block quadratics but is not universal. On (H) and (R) with $m=256$
(Section~\ref{sec:exp-A4}), L-Katyusha without its negative momentum diverges for every $\kappa\ge10^3$ ($n=16$), and
Nesterov's momentum on the SVRG estimator, even re-tuned at every $\kappa$ over six settings, has the log--log slope
$0.81$ on both instances instead of $\frac12$ ($95\%$ intervals $[0.62,0.99]$ and $[0.65,0.98]$, wide because its
profile is irregular) and needs $3.0$ times the gradients of L-Katyusha at $\kappa=10^4$. In Section~\ref{sec:exp-A2}
(one seed, $\kappa=10^4$), the best of fifteen settings needs $2.9$ times the gradients of L-Katyusha to reach $10^{-6}$
on (H) with $m=64$, a chain too short for the lower bound to be certified, whereas on logistic regression both
anchor-free schemes converge and need about half the gradients of L-Katyusha with its worst-case parameters. The
proof-level obstruction of \resref{prop:naive-kappa} is a limit of the argument, not a lower bound, as
\resref{rem:naive-honest} details, and our evidence comes from small grids.

\paragraph{A4: exponents on hard instances.}
Confirmed inside the window of $\kappa$ where \resref{prop:A4} certifies the exponents, with two corrections to the
expectations of the proposal. On four quadratic instances ($n=16$, 2 or 3 seeds) the log--log slopes of the gradient
counts over $\kappa\ge10n$ are $0.95$--$1.02$ for SAGA ($0.99$ inside the window on (S)), $1.00$--$1.02$ for L-SVRG (run on (H)
and (R) only) and $0.49$--$0.52$ for L-Katyusha. First, on block-separable instances GD and AGD share these exponents
with smaller constants, because $L_F\approx R^2/n$; the orders $n\kappa$ and $n\sqrt\kappa$ of the proposal appear only
on the aligned instance, where GD costs $14$--$16$ times and AGD $3.9$--$4.4$ times more than on the block instance (S).
Second, the exponents are tied to a finite window: beyond $\kappa_{\max}$ the chain of (S) is exhausted
($\kappa_{\rm eff}/\kappa=0.44$ at the nominal $\kappa=10^6$), and the slopes of SAGA and GD against the nominal
$\kappa$ drop from $0.99$ inside the window to $0.95$ and $0.94$ ($95\%$ intervals $[0.90,0.99]$ for both; $1.03$ against
$\kappa_{\rm eff}$).

\needspace{5\baselineskip}% keeps the run-in heading with the list (needspace is loaded in main.tex)
\paragraph{Limitations.}
\begin{itemize}[leftmargin=1.4em,itemsep=1pt,topsep=2pt]
\item \emph{Data.} All logistic problems are synthetic: Gaussian features with power-law spectra, $n=1000$, $d=50$,
one data set per $\alpha$ shared by all $\kappa$ (which only set $\lambda$), $\theta_0=0$, and isotropic features with
$d=40$ and $n\in\{10^3,10^4\}$ for the validation of Section~\ref{sec:exp-validation}. Real data, other ratios $d/n$ and
nearly separable data were not tested, and the reproduction of the experiment of the book is only qualitative.
\item \emph{Seeds and intervals.} We used 5 seeds (Sections~\ref{sec:exp-B3} and~\ref{sec:exp-B1}), 3
(Figure~\ref{fig:C}, Section~\ref{sec:exp-A3}, (H) and (R) in Section~\ref{sec:exp-A4}), 2 ((S) and (A)) or 1
(Sections~\ref{sec:exp-validation}, \ref{sec:exp-B2} and~\ref{sec:exp-A2}); bands are min--max ranges over the seeds,
not confidence intervals. For the stochastic methods the seeds are independent replicates (on (R) each seed also draws
its own instance), whereas GD and AGD are deterministic, so that on (H), (S) and (A) their counts coincide for every
seed. The slopes are least-squares fits of the per-$\kappa$ geometric means over the seeds, and their $95\%$
$t$-intervals ($K-2$ degrees of freedom, $K$ values of $\kappa$) measure the departure of these means from a single
power law over a finite window of $\kappa$, the dominant uncertainty on an exponent, rather than sampling noise:
additional seeds would not reduce it, no coverage is claimed, and an exponent remains a local, window-dependent
quantity.
\item \emph{Attribution.} The shares depend on the order in which the inequalities are undone, and the two orders of
\resref{prop:B3-attribution} differ in the memory regime; the shares are additive in contraction exponents and not
multiplicative factors, and the product decomposition used for $\kappa\ge10^4$ is approximate ($91$--$96\%$ of the
speed-ups).
\item \emph{Exactness.} Conditional expectations are exact only along the realized trajectories; products of
conditional factors are realized quantities, and the trajectory-dependent bounds (\resref*{prop:B3-traj}) are
a-posteriori, not guarantees.
\item \emph{Precision floors.} Nothing is analysed below $V_t=10^{-18}$ for SAGA ($\norm{\theta_t-\etastar}^2=10^{-18}$
in Section~\ref{sec:exp-B2}) or $\Psi_t=10^{-10}$ for L-Katyusha, so the late phase of the well-conditioned runs is
truncated; conversely, the runs with $\kappa=10^5$ and $\alpha=2$ do not reach the neighbourhood of $\etastar$ in 100
passes.
\item \emph{Hard instances.} One $n$ ($16$), one accuracy ($\varepsilon=10^{-8}$) and a finite window of $\kappa$; the
dependence on $n$ and the regime $\kappa\ll n$ were not explored, and the instance (H) with $m=64$ used in
Section~\ref{sec:exp-A2} is too short for the lower bound to be certified at the accuracies used there.
\item \emph{Tuning and counting.} The naive scheme was tuned on small grids (fifteen settings in
Section~\ref{sec:exp-A2}; six per $\kappa$ in Section~\ref{sec:exp-A4}, selected on seed $0$ on (H) and on each seed on
(R)), while L-Katyusha kept its theoretical parameters; a finer grid, or one in which $\beta$ and $\gamma$ scale with
$\kappa$ beyond the single rule $\beta=(\sqrt\kappa-1)/(\sqrt\kappa+1)$, could change the measured exponent of the
naive scheme. The reference gradients $\fip(w)$ of the SVRG-type methods are charged as if recomputed at every
iteration (the code stores them, as SAGA stores its table), which raises their counts by about one half compared with
the convention with storage used by the bounds of Appendix~\ref{sec:lower_bound} ($3$ instead of $2$ gradients per
iteration on average); this changes constants, not exponents, but it decides the order of SAGA and L-Katyusha on
logistic regression in Section~\ref{sec:exp-A2} (with storage L-Katyusha needs $7.1\cdot10^4$ gradients to
$10^{-10}$, against $8.8\cdot10^4$ for SAGA).
\end{itemize}

\paragraph{Open questions.}
(i) The exact worst-case one-step contraction of the Lyapunov function of Proposition~\ref{prop:59}, computable by
performance estimation when interpolation conditions are known for its class of
functions~\cite{taylor2017smooth,goujaud2024pepit}
and bounded from above by the linear matrix inequality of~\cite{hu2017unified}, which only uses quadratic constraints
between the iterate and $\etastar$ (see also~\cite{taylor2019stochastic}), would separate the slack of the proof from
the benign character of our instances, and its dual multipliers would say which inequalities are active in the worst
case, to be compared with our trajectory-wise attribution. (ii) An a-priori version of the trajectory-dependent bound
can be obtained on a region (\resref*{prop:B3-region}), with the local strong convexity of logistic regression on a
ball around $\etastar$ (\resref*{prop:B3-logistic}); how much of the gap such a computable certificate recovers is
open. (iii) Once the curvature is handled, the remaining slack of Proposition~\ref{prop:59} is in its memory
coefficient and in the step I5 (Section~\ref{sec:exp-B1}). At the textbook step a larger weight $\Delta$ already
certifies up to twice the memory rate $\frac1{3n}$ when $n\gg\kappa$, by \resref{cor:B3-c}, and \resref{prop:B3-optimum}
allows at most a factor $3$ over all $(\gamma,\Delta)$. For $\kappa\le\frac{9n}{16}$, where the textbook rate is
$\frac1{3n}$, this factor $3$ is also the most that any analysis of $\E[V_t]$, whatever $\Delta$, can gain, since, by
\resref{rem:B3-conclusion}, no bound on $\E[V_t]$ certifies more than $\frac1n$ per iteration: the open question is
whether another Lyapunov function certifies more at the textbook step. (iv) The alignment of SAGA with the flattest
direction for
$\kappa\gg n$, seen in two runs, should be tested on more configurations, for instance by recording the spectrum of the
noise of the estimator. (v) No lower bound shows that Nesterov's momentum on the SVRG estimator with one sample per
iteration cannot accelerate; measuring its exponent as a function of the mini-batch size would connect our
observations with the mini-batch analyses of~\cite{nitanda2016accelerated,driggs2022accelerating}. (vi) The largest
additive share for $\kappa\ge10^4$ belongs to cocoercivity with $\max_iL_i$, which non-uniform sampling replaces by an
average of the $L_i$~\cite{xiao2014proximal,schmidt2015nonuniform,needell2016stochastic}, removing the part of this
slack due to the spread of the $L_i$; instrumenting such a variant, and exploiting $\muEff$ in accelerated methods
through restarts~\cite{fercoq2019adaptive,roulet2020sharpness} within the limits set
by~\cite{guilleescuret2022gradient,arjevani2017limitations}, are natural next steps.
